% !TEX root = ../main.tex
\section{Related Work}
\label{sec:related}

This paper draws on nine bodies of work and departs from each in the same way: it
treats the object of interest as institutional reliance on a \emph{particular}
assertion rather than as an aggregate property of a model or a distribution. The
contribution is integration and operationalisation, not discovery in any of these
fields, and the sections below say specifically what is taken and where the departure
lies.

\paragraph{Attribution, faithfulness, and citation verifiability.} The problem of
determining whether generated text is attributable to identified sources has a
careful operational treatment, and the framework of \citet{rashkin2023} is what the
audit procedure of Section~\ref{sec:measure-rate} adopts. The distinction between
consistency with a source and correspondence with the world originates in work on
summarisation \citep{maynez2020} and is established as standard by surveys of the area
\citep{ji2023}; this paper claims no priority for it. \citet{liu2023} measure how
often statements in generative search output are actually supported by the citations
attached to them, which is the closest existing measurement to the one proposed here.

The departure is that attribution is necessary and not sufficient. A citation can be
present, accurate, and authoritative while still failing temporal validity,
applicability to $c$, standing for the decision class, or authorization for reliance.
Attribution establishes I12. Warrant requires I8--I19.

\paragraph{Provenance, lineage, and temporal data management.} The requirement that a
derived value carry the record of what produced it is not new; it is the subject of a
substantial database literature, surveyed by \citet{cheney2009} under the
why/how/where taxonomy of provenance. The valid-time and transaction-time distinction
this paper leans on in Section~\ref{sec:resolved-context} is standard bitemporal
modelling \citep{snodgrass1999}, and none of it is offered as novel.

What this paper does is apply those requirements to machine-generated assertion, where
they are routinely absent. The observation that an enterprise would not accept a
warehouse fact without lineage, and does accept a model assertion without it, is the
whole of the argument in Section~\ref{sec:canonical-form}.

\paragraph{Retrieval-augmented generation and grounded generation.} Retrieval
conditioning \citep{lewis2020} is the architectural ancestor of
Section~\ref{sec:governed-rag}, and it does improve factuality and supply provenance
of a kind. Section~\ref{sec:retrieval-insufficient} states the departure precisely
rather than dismissively: standard retrieval supplies authenticity where the corpus is
controlled and traceability where the event is logged, and supplies none of the four
conditions that are relations to $c$ or to the evidential base as a whole.
% TODO(cite): metadata-filtered and structured/hybrid retrieval, knowledge-graph
% grounding, and neuro-symbolic or rule-based eligibility systems. Several of the
% controls in Section 11 are standard practice in that literature and it should be
% credited rather than implied to be new here.

\paragraph{Uncertainty, selective prediction, and abstention.} A substantial
literature addresses whether models know what they know and how uncertainty can be
surfaced \citep{kadavath2022,lin2022}, and conformal methods give distribution-free
coverage guarantees under stated exchangeability assumptions
\citep{angelopoulos2023}. The abstention design of Section~\ref{sec:abstention}
borrows the risk-coverage framing directly.

The departure is the one recorded at I5 and generalised in
Section~\ref{sec:independence}: these are properties of a distribution, and reliance
decisions are made about individual assertions. A calibrated probability of
correctness is not an audit trail, does not identify the evidence relied upon, and
does not establish applicability to the case at hand. Conformal coverage is the
closest thing to a per-instance guarantee in this family, and it guarantees coverage
of a label set, not that the asserted content rests on evidence that governs the
decision.
% TODO(cite): selective classification / learning-to-defer. I could not verify a
% specific reference for the risk-coverage tradeoff in deep networks; supply one.

\paragraph{Robustness, paraphrase consistency, and behavioural testing.} The
observation that model behaviour moves under semantically inert input variation is
well established: format and prompt-design sensitivity \citep{sclar2024},
paraphrase-consistency measurement \citep{elazar2021}, insensitivity to prompt meaning
\citep{webson2022}, and framing effects under user preference cues
\citep{sharma2024}. Behavioural and capability-based testing beyond aggregate accuracy
is the programme of \citet{ribeiro2020}, and the form-invariance audit of
Section~\ref{sec:form-audit} is a specialisation of it.

The departure is what the instability is taken to \emph{mean}. That literature treats
paraphrase sensitivity primarily as a robustness defect to be reduced. This paper
treats it as the reason a validation inference fails, and therefore as an argument for
measuring a different property rather than for hardening the same one.

\paragraph{Assurance cases and model risk management.} The argument here belongs to a
tradition older than the technology: structuring a claim, the evidence for it, the
argument connecting them, the controls applied, and the accountable party.
Safety-case and assurance-case practice in regulated engineering supplies the shape of
Table~\ref{tab:invariants-system}, and independent model validation as a second-line
function supplies the fourth row of Table~\ref{tab:three-teams}.
% TODO(cite): structured assurance-case notation, safety-case practice, and
% supervisory model-risk guidance. I could not verify specific references and have
% deliberately not guessed; the author is better placed to cite the governing
% standards for the target sector.

\paragraph{Legal, temporal, and jurisdiction-aware reasoning.} The phenomena this
paper indexes into $c$ are each recognised in their own literatures --- temporal
validity in time-sensitive question answering, jurisdiction in legal information
retrieval, population scope in clinical decision support.
% TODO(cite): time-sensitive QA, legal-domain retrieval and reasoning, and
% jurisdiction-aware systems. This is the literature most likely to contain a prior
% statement of the applicability condition, and the priority claim in
% Section~\ref{sec:novelty} should not be made until it has been checked.
What appears to be absent, and what Section~\ref{sec:novelty} claims, is their
integration into a single condition on grounded assertion together with the
observation that faithfulness metrics cannot express it.

\paragraph{Human factors and automation bias.} That operators over-trust automated
output, and under-verify it when it appears authoritative, is long established
\citep{parasuraman1997,skitka1999}. This bears directly on the argument of
Section~\ref{sec:cell-ii}: a fluent assertion with an accurate citation is precisely
the presentation most likely to suppress the verification it requires, which is why
cell (ii) is dangerous rather than merely mislabelled. It also bears on
Section~\ref{sec:relocation}: rendering a resolved request back for confirmation is a
control that automation-bias research predicts will be under-used, and it should be
designed accordingly.

\paragraph{Causal attribution and the limits of ablation.} The
sensitivity condition of Definition~\ref{def:grounding} is a counterfactual claim
about the production of an assertion, and estimating it inherits every difficulty of
attribution in generative systems. Section~\ref{sec:sensitivity-protocol} sets out the
controlled design this paper requires and is explicit that the protocol
\emph{estimates} rather than establishes dependence.
% TODO(cite): RAG attribution methods and known limitations of ablation-based
% attribution. The confound identified in Section~\ref{sec:sensitivity-protocol} is
% likely to be documented in that literature and should be credited if so.

\paragraph{Alignment and output shaping.} Preference optimisation
\citep{ouyang2022} and rule-conditioned training \citep{bai2022} shape output
distributions and improve compliance behaviour; Section~\ref{sec:rlhf} treats the
first as part of the object of analysis rather than as a competing proposal. The
departure follows from the structure of \eqref{eq:rlhf}: the reward is a function of
the output, so unless the training and deployment architecture explicitly represents
and verifies the evidence relation, these methods cannot establish a source-specific,
version-specific, context-specific evidence path for a live assertion.
